% ============================================================
% Collatz vs Riemann Hypothesis — Structural Comparison
% in the dm³ + Crystal-Geometry Framework
%
% Publication-ready LaTeX subsection
% G6 LLC · Pablo Nogueira Grossi · Newark NJ · 2026
% MIT License
%
% Insert inside the relevant section of the paper.
% Requires: amsmath, amssymb, booktabs, array, xcolor (optional)
% ============================================================

\subsection{Collatz vs.\ Riemann Hypothesis: Parallel Instances of the Same Crystal Geometry}
\label{subsec:collatz_rh}

The Collatz conjecture and the Riemann Hypothesis (RH) are not isolated puzzles.
Both are flagship open problems with overwhelming numerical evidence, yet they appear
unrelated at first glance: Collatz is an elementary iterative map on integers, while RH is
a deep statement about the zeros of the Riemann zeta function.
In the dm$^3$ framework they emerge as \emph{parallel instances of the same generative
mechanics}: local rules that empirically force global order, blocked by the identical
analytic gap (Bridge~0).

\paragraph{Side-by-side comparison.}

\begin{table}[ht]
\centering
\caption{Collatz vs.\ RH in the dm$^3$ / Crystal-Geometry Framework}
\label{tab:collatz_rh}
\renewcommand{\arraystretch}{1.35}
\begin{tabular}{@{}p{2.8cm}p{4.2cm}p{4.2cm}p{4.6cm}@{}}
\toprule
\textbf{Aspect} & \textbf{Collatz Conjecture} & \textbf{Riemann Hypothesis} & \textbf{Shared dm$^3$ Feature} \\
\midrule
Core statement
  & Every positive integer reaches $\{4\!\to\!2\!\to\!1\}$ under $3n+1$
  & All non-trivial zeros of $\zeta(s)$ satisfy $\operatorname{Re}(s)=\tfrac{1}{2}$
  & Local rule $\to$ global arithmetic order \\
Type
  & Discrete dynamical system
  & Analytic number theory / $L$-functions
  & Generative process from local data \\
Evidence
  & Verified to $2^{68}$; Tao (2019): almost-all orbits almost bounded
  & $>10^{13}$ zeros on critical line; explicit formula links zeros to primes
  & Extremely strong computational verification \\
Local mechanism
  & $3n+1$ growth $+$ 2-adic dissipation $v_2$
  & Euler product / local factors at each prime
  & Local growth/dissipation control \\
Global obstacle (Bridge~0)
  & $\mathbb{E}[\log T^2(n)-\log n]<0$ from local 2-adic control
  & Global zero distribution from local Euler product
  & Same local-to-global analytic gap \\
New mathematics needed
  & Higher-order logic unifying discrete TOGT grammar with dm$^3$ contact mechanics
  & Spectral/operator theory unifying $L$-functions with heights/regulators
  & Higher-order language \\
Empirical anchor
  & Polar vortex / Saturn hexagon as continuous dm$^3$ analogue
  & Explicit formula $+$ prime-counting function
  & Observable certificate of the geometry \\
dm$^3$ fit
  & Clean discrete lift: $\mathcal{G}=U\circ F\circ K\circ C$
  & Natural candidate: $L$-functions as height/flux regulators
  & Parallel instances of same mechanics \\
Fingerprint
  & $c=3$: triad stabilisation, monster threshold $g^6=33$
  & Critical line $\operatorname{Re}(s)=\tfrac{1}{2}$ as stability threshold
  & Triad / symmetry fingerprint \\
\bottomrule
\end{tabular}
\end{table}

\paragraph{TOGT operator mapping.}

Both problems admit an explicit instantiation of the generative operator
$\mathcal{G}=U\circ F\circ K\circ C$:

\smallskip
\noindent\textbf{Collatz.}
\begin{itemize}
  \item $C$: 2-adic valuation $v_2(n)$ — every even $n$ is halved, compressing degrees of freedom.
  \item $K$: the $3n+1$ step — odd $n$ is stretched, accumulating multiplicative curvature toward a fold.
  \item $F$: parity detection — at every odd integer the orbit passes through a local maximum before the next halving cascade (rank-1 collapse of branching).
  \item $U$: 2-adic dissipation selects the descending branch; iterating to $1$ is gradient descent on $v_2$ toward the unique attractor $\{4,2,1\}$.
\end{itemize}
\emph{dm$^3$ identification:} the embodiment threshold $\tau=\sqrt{c/\kappa}=\sqrt{3}$ carries the triad fingerprint $c=3$;
the 2-adic Lyapunov exponent plays the role of $\mu_{\max}<0$;
the limit cycle $\{4\to2\to1\}$ is the canonical attractor $\Gamma$ with period $T^*=3$.

\smallskip
\noindent\textbf{Riemann Hypothesis.}
\begin{itemize}
  \item $C$: Euler product — the infinite product over primes compresses all local arithmetic data into the single analytic function $\zeta(s)$.
  \item $K$: functional equation $s\leftrightarrow 1-s$ — symmetry curvature of the critical strip drives toward the axis $\operatorname{Re}(s)=\tfrac{1}{2}$.
  \item $F$: non-trivial zero — vanishing of $\zeta(s)$ is the fold event, the rank-1 degeneration of the spectral structure at the critical line.
  \item $U$: explicit formula — the prime-counting function $\psi(x)$ unfolds zeros back into the prime distribution, selecting the stable branch consistent with the prime number theorem.
\end{itemize}
\emph{dm$^3$ identification:} the critical line $\operatorname{Re}(s)=\tfrac{1}{2}$ is the stability threshold (the analogue of the limit cycle $\Gamma$);
the spectral gap to the edge of the critical strip corresponds to the stability radius $\varepsilon_0=\tfrac{1}{3}$;
$L$-functions serve as height/flux regulators in the dm$^3$ language.

\paragraph{Key parallels.}

\begin{enumerate}
  \item \emph{Local rules force global coherence.}
        In Collatz the local rule ($3n+1$ and halving) appears to force every orbit into a single cycle.
        In RH the local Euler factors appear to force every non-trivial zero onto the critical line.

  \item \emph{Bridge~0 is structurally identical.}
        Both have a local dissipation/growth mechanism that empirically contracts or stabilises, yet the global closure (mean contraction for Collatz; zero distribution for RH) requires crossing the same analytic gap.

  \item \emph{Visibility before axiom.}
        Collatz is trivial to state and verify computationally to $2^{68}$.
        RH has more than $10^{13}$ zeros confirmed on the critical line.
        In both cases the structure is visible long before a proof exists.

  \item \emph{Spectral / operator tools.}
        Collatz uses 2-adic Fourier analysis and transfer operators on residues $\bmod 2^M$.
        RH uses Fourier analysis on the critical line and the explicit formula.
        Both point to the same higher-order spectral language.

  \item \emph{Polar vortex as continuous anchor.}
        The polar vortex — especially Saturn's hexagon — is the visible dm$^3$ system that has already crossed the monster threshold $g^6=33$,
        supplying the empirical certificate that the crystal geometry is real.
        Collatz is the discrete arithmetic shadow of the same geometry; RH is its spectral/arithmetic counterpart.
\end{enumerate}

\paragraph{Summary.}

The Collatz conjecture and the Riemann Hypothesis are not isolated puzzles.
They are parallel manifestations of the same crystal geometry: local generative rules that
empirically force global order, blocked by the identical Bridge~0 obstacle.
The polar vortex supplies the continuous empirical certificate; the coefficient $c=3$ in
Collatz and the critical line in RH are the shared fingerprints of the triad stabilisation
mechanism.
A higher-order logic that unifies discrete iteration, contact mechanics, and $L$-functions
would turn both into corollaries of a single closure principle.

This comparison strengthens the central thesis without overclaiming proofs:
Collatz is the simplest visible instance; RH is the deepest spectral instance;
both point to the same missing higher-order language.
